\begin{tikzpicture}[
  font=\small,
  node distance=5mm and 6mm,
  box/.style={draw, rounded corners=2pt, align=center, minimum height=11mm, minimum width=23mm, fill=gray!10},
  module/.style={draw, rounded corners=2pt, align=center, minimum height=13mm, minimum width=31mm, fill=green!10},
  train/.style={draw, dashed, rounded corners=2pt, align=center, minimum height=10mm, minimum width=28mm, fill=orange!12},
  arr/.style={-{Latex[length=2mm]}, thick}
]
\node[box] (cams) {左前 / 正前 / 右前\\RGB 图像};
\node[module, right=of cams] (fpn) {预训练 ResNet-101\\+ 多尺度 FPN};
\node[module, right=of fpn] (bev) {GVAD 增强的\\6 层 BEV 编码器};
\node[module, right=of bev] (decoder) {3 层占用解码器\\$100\!\times\!100\!\times\!25\!\times\!6$};
\node[box, right=of decoder] (out) {当前三维\\语义占用};
\draw[arr] (cams)--(fpn);
\draw[arr] (fpn)--(bev);
\draw[arr] (bev)--(decoder);
\draw[arr] (decoder)--(out);
\node[box, below=8mm of bev] (nearfar) {近场稠密 + 远场稀疏\\恢复后最终稠密交叉注意力};
\draw[arr] (nearfar)--(bev);
\node[train, below=8mm of decoder] (adhr) {训练期 ADHR：不确定体素 +\\作物/自由空间困难间隙};
\draw[arr] (decoder)--(adhr);
\node[align=left, below=1mm of cams, text width=43mm] {\footnotesize GVAD：局部可变形采样 + 基于相机投影有效掩码的 32 个可见锚点。};
\end{tikzpicture}

